% !TEX root = ../../hw01.tex
% Source: ma-gy7033fall2026homework1.pdf, p. 3, Problem 9(a--c), including hint.
% Omit only "from the lecture" in the hint for self-contained shared wording.
% Lucas authorized adding the prerequisite comparison theorem on November 9.
\begin{exercise}[Extending an independent set]
  \label{ex:hw01:basis-extension}
  Let \(V\) be a vector space over \(\mathbb F\), where
  \(\mathbb F\) denotes \(\R\) or \(\C\), and let \(V\) have dimension \(n\).
  \begin{enumerate}[(a)][2]
    \item Prove that every linearly independent subset of \(V\) is finite
      and can be extended to a basis of \(V\).
    \item Deduce that any \(n\) linearly independent vectors form a basis,
      and that any set of \(n\) vectors spanning \(V\) is a basis.
    \item Prove that a subspace \(W\subseteq V\) with \(\dim W=n\)
      must equal \(V\). Give an example showing that an infinite dimensional
      space can have a proper subspace of the same dimension; exhibit bases
      to justify the comparison.
  \end{enumerate}
  \emph{Hint.} For (a), show that adjoining \(v\notin\operatorname{span}(S)\)
  to an independent set \(S\) preserves independence. Use the comparison
  theorem to bound the number of steps.
  \solutionlink{hw01:basis-extension}
\end{exercise}
